\documentclass[a4paper,11pt]{article}

% ---------- Кодировка и язык ----------
\usepackage[utf8]{inputenc}
\usepackage[T2A]{fontenc}
\usepackage[english,russian]{babel}

% ---------- Геометрия страницы ----------
\usepackage[margin=2.5cm]{geometry}

% ---------- Математика ----------
\usepackage{amsmath,amssymb,amsthm,mathtools}
\usepackage{bm}          % жирные символы в формулах: \bm{x}

% ---------- Картинки ----------
\usepackage{graphicx}
\graphicspath{{figures/}}   % картинки лежат в latex/figures/
\usepackage{float}          % позиция [H] для картинок
\usepackage{subcaption}     % несколько картинок рядом

% ---------- Оформление ----------
\usepackage{enumitem}
\usepackage{booktabs}       % красивые таблицы
\usepackage{xcolor}
\usepackage[colorlinks=true,linkcolor=blue,urlcolor=blue,citecolor=blue]{hyperref}
\usepackage{cleveref}       % \cref{...} для умных ссылок
\newcommand{\lean}[1]{\mbox{\path{#1}}} % имена из Lean-кода, без переносов внутри имени (пакет url)

% ---------- Теоремы и окружения ----------
\theoremstyle{plain}
\newtheorem{theorem}{Теорема}[section]
\newtheorem{lemma}[theorem]{Лемма}
\newtheorem{proposition}[theorem]{Утверждение}
\newtheorem{corollary}[theorem]{Следствие}

\theoremstyle{definition}
\newtheorem{definition}[theorem]{Определение}
\newtheorem{example}[theorem]{Пример}

\theoremstyle{remark}
\newtheorem*{remark}{Замечание}

% ---------- Полезные макросы ----------
\newcommand{\R}{\mathbb{R}}
\newcommand{\Z}{\mathbb{Z}}
\newcommand{\N}{\mathbb{N}}
\newcommand{\Cx}{\mathbb{C}}   % \C занят babel-russian
\newcommand{\norm}[1]{\left\lVert #1 \right\rVert}
\newcommand{\abs}[1]{\left\lvert #1 \right\rvert}
\newcommand{\set}[1]{\left\{ #1 \right\}}
\newcommand{\ip}[2]{\langle #1, #2\rangle} % скалярное произведение
\newcommand{\Rpp}[1][2]{\R^{#1}_{++}} % \Rpp = R^2_{++}, \Rpp[n] = R^n_{++}
\newcommand{\Rp}[1][2]{\R^{#1}_{+}}   % \Rp = R^2_{+}, \Rp[n] = R^n_{+}

% Заметка на полях / TODO
\DeclareMathOperator{\conv}{conv}
\DeclareMathOperator{\cone}{cone}
\DeclareMathOperator{\Sp}{Sp}
\DeclareMathOperator{\co}{co}
\DeclareMathOperator*{\argmax}{arg\,max}
\usepackage{array}
\usepackage{empheq}
\usepackage{xurl}
\newcommand{\todo}[1]{\textcolor{red}{\textbf{TODO:} #1}}

\title{Альтернативная ветка: тесты рациональности\\ в теории выявленных предпочтений}
\author{Обзор моделей для сравнения с заметками о неоклассической разрешимости}
\date{8 октября 2026}

\begin{document}
\maketitle
\tableofcontents

\section*{Зачем этот обзор и что прочитано}

В заметках доказано: вектор выпусков $y$ неоклассически разрешим при ценах $P$ тогда и только тогда, когда $y\in\cone\set{\mathbf 1_S\mid S\in\Sp(P)}$, где $\Sp(P)$~--- достижимые спектры, описанные через покрытие. Эта линия идёт от модели Хаутеккера--Йохансена через Шананина, Хенкина, Молчанова и Агальцова. Независимо в западной эконометрике сложилась ветка с той же математической формой «наблюдение лежит в конусе 0/1-векторов допустимых типов», но про потребителей. Ниже каждая модель этой ветки описана по первоисточнику: примитивы, данные, гипотеза, теорема, алгоритм.

Прочитаны полностью: Kitamura--Stoye, рабочая версия 2012 (разделы 3--8); Stoye 2019 (ARSP); Smeulders 2018 (генерация столбцов); Koida--Shirai 2024 (двойственная характеристика); Deb--Kitamura--Quah--Stoye 2023; Dean--Martin; Demuynck--Rehbeck 2021; Varian 1990; Aguiar--Kashaev 2020. Формулы ниже транскрибированы из этих текстов с сохранением нумерации, где она есть. Статья Varian 1985 недоступна, о ней только по пересказу в Varian 1990.

\section{Детерминированная теория: один потребитель}

\subsection{Данные, выявленное предпочтение, GARP, теорема Африата}

Данные~--- $T$ наблюдений $(p^t,x^t)$, $p^t\in\Rpp[K]$ цены, $x^t\in\Rp[K]$ выбранный набор. \textit{Прямое выявленное предпочтение}: $x^tR^Dx^s$, если $p^tx^t\ge p^tx^s$, то есть $x^s$ был доступен, когда выбрали $x^t$. Строгое: $x^tP^Dx^s$, если $p^tx^t>p^tx^s$. $R$~--- транзитивное замыкание $R^D$.

\textbf{GARP.} Если $x^tRx^s$, то $p^sx^s\le p^sx^t$. Иначе: нет цепочки $x^{t_1}R^D\cdots R^Dx^{t_M}$ с $x^{t_M}P^Dx^{t_1}$.

\textbf{Теорема Африата (1967).} Данные рационализуемы локально ненасыщаемой функцией полезности ($p^tx^t\ge p^tx\Rightarrow u(x^t)\ge u(x)$) тогда и только тогда, когда они удовлетворяют GARP; при этом $u$ можно взять непрерывной, монотонной и вогнутой.

Эквивалентная форма через «числа полезности» (Demuynck--Rehbeck, теорема~2, через расширение Шпильрайна): GARP выполнен тогда и только тогда, когда существуют $u_t\in[0,1]$ с
\[
    p^tx^t\ge p^tx^v\Rightarrow u_t\ge u_v,\qquad p^tx^t>p^tx^v\Rightarrow u_t>u_v.
\]
Эта форма~--- рычаг для всех MILP ниже: GARP превращается в существование чисел, удовлетворяющих линейным ограничениям с индикаторами.

\subsection{Вариант в пространстве цен: Deb--Kitamura--Quah--Stoye}

Это ближайший к нам детерминированный объект: отношение предпочтения задано на \emph{ценах}.

\textbf{Модель.} Потребитель покупает наблюдаемые $L$ благ, не обязательно все свои блага. Функция полезности с учётом расходов $U(x,-e)$, $U:\Rp[L]\times\R_-\to\R$, строго убывает по расходу $e$. При ценах $p$ решается
\[
    V(p):=\sup_{x\in\Rp[L]}U(x,-p\cdot x).
\]
Данные $\set{(p^t,x^t)}$ \textit{рационализуемы}, если $x^t\in\argmax_xU(x,-p^t\cdot x)$ при всех $t$. Это не рационализуемость по Африату: бюджетного ограничения нет, цена набора входит в полезность.

\textbf{Выявленное предпочтение на ценах.} $p^{t'}\succeq_pp^t$ ($\succ_p$), если $p^{t'}\cdot x^t\le p^t\cdot x^t$ (строго): набор $x^t$ при ценах $p^{t'}$ обошёлся бы не дороже. $\succeq_p^*$, $\succ_p^*$~--- транзитивные замыкания.

\textbf{GAPP.} Нет $t,t'$ с $p^{t'}\succeq_p^*p^t$ и $p^t\succ_p^*p^{t'}$.

\textbf{Теорема 1 (DKQS).} Равносильны: (1) данные рационализуемы функцией $U$; (2) GAPP; (3) рационализуемы строго возрастающей, непрерывной, вогнутой $U$, у которой максимум достигается при всех $p$.
Доказательство сводится к Африату: добавляется $(L+1)$-е благо с ценой $1$ и спросом $M-p^t\cdot x^t$, и GAPP для исходных данных равносилен GARP для расширенных.

GAPP и GARP логически независимы, но совпадают на \textit{изорасходных} данных $\check x^t:=x^t/(p^t\cdot x^t)$ (предложение~1). Это понадобится в стохастической версии.

\subsection{Индексы согласия}

Когда GARP нарушен, спрашивают, насколько. Все индексы ниже относятся к одному потребителю.

\textbf{Индекс Африата (CCEI).} Ослабленное отношение $x^tR^D(e)x^s\iff e\,p^tx^t\ge p^tx^s$, $e\in[0,1]$. $\mathrm{AEI}=\max\set{e\mid R(e)\text{ ациклично}}$. Один номер на весь набор данных, измеряет худший цикл, вычислим за полиномиальное время бинарным поиском.

\textbf{Вектор Варьяна (1990).} Обобщённая функция компенсации $m(p,x,S)=\inf_{ySx}py$ и поэлементный индекс
\[
    i^t=\frac{m(p^t,x^t,R)}{p^tx^t},
\]
равный $1$ при выполнении GARP и $p^tx^s/p^tx^t$ при нарушении парой $(s,t)$. Предложение~1: данные удовлетворяют $\mathrm{GARP}(i^t)$. Вектор $(i^t)$ не минимален: все наблюдения на одном цикле получают $i^t<1$, хотя достаточно возмутить одно. Уточнённый вектор $v^t$ строится итеративно: найти циклы, разорвать каждый в наблюдении с максимальным нетривиальным $i^t$, пересчитать, повторять до выполнения GARP. $v^*=\min_tv^t$~--- нижняя оценка индекса Африата. Варьян (1985), по его же пересказу,~--- настоящее минимальное возмущение количеств, но трудоёмкое.

\textbf{Houtman--Maks.} $\mathrm{HMI}=\max\set{|A|/T\mid (p^t,x^t)_{t\in A}\text{ удовлетворяет GARP}}$. Только число нарушений, без их величины. NP-трудно.

\textbf{Minimum Cost Index (Dean--Martin).} Стоимость отношения $x^tR^Dx^s$~--- денежная метрика $p^t(x^t-x^s)\ge0$. Индекс
\[
    \mathrm{MCI}=\min_{B\subset R^D}\frac{\sum_{(t,s)\in B}p^t(x^t-x^s)}{\sum_tp^tx^t}\quad\text{при условии, что }R^D\setminus B\text{ ациклично}.
\]
Это задача о минимальном по весу разрывающем множестве дуг. NP-трудно; авторы сводят к покрытию множеств и решают стандартными солверами. $\mathrm{MCI}=0$ тогда и только тогда, когда выполнен SARP. Сравнение индексов по двум осям, «серьёзность нарушения» и «агрегирование»: AEI~--- слабейшее звено и максимум по циклам; MCI~--- слабейшее звено и сумма; HMI~--- без веса и сумма; Money Pump Index (Echenique--Lee--Shum)~--- сумма звеньев цикла и среднее по циклам.

\subsection{MILP для индексов: Demuynck--Rehbeck}

Все индексы выше записываются как смешанно-целочисленные линейные задачи с одной и той же конструкцией.

\textbf{Кодирование GARP.} Бинарные $U_{t,v}\in\set{0,1}$ означают «$u_t\ge u_v$». Фиксируем $0<\varepsilon<1/T$, $\alpha>\max_tp^tx^t$ и $\delta>0$ меньше всех положительных разностей $p^tx^v-p^tx^t$:
\begin{align*}
    u_t-u_v&\le-\varepsilon+2U_{t,v}, & U_{t,v}-1&\le u_t-u_v, \tag{IP-1,2}\\
    p^tx^t-p^tx^v&\le-\delta+\alpha U_{t,v}, & \alpha(U_{v,t}-1)&\le p^tx^v-p^tx^t. \tag{IP-3,4}
\end{align*}
(IP-1,2) связывают $U$ с числами $u$, (IP-3,4)~--- с данными. Следствие~1: GARP выполнен тогда и только тогда, когда система совместна.

\textbf{HMI.} Бинарные $A_t$ (наблюдение оставлено), в (IP-3,4) добавляется $\alpha(1-A_t)$, цель $\max\sum_tA_t/T$.

\textbf{Средний индекс Варьяна.} Переменные $e_t\in[0,1]$ входят в (IP-3,4) линейно: $e_tp^tx^t-p^tx^v\le\alpha U_{t,v}$. Цель $\max\sum_te_t/T$, либо квадратичная $\min\sum_t(1-e_t)^2/T$.

\textbf{MCI.} Бинарные $B_{t,v}$ (ребро удалено), $\alpha(U_{t,v}+B_{t,v})$ в правой части, цель
\[
    \min\ \sum_{t,v}B_{t,v}\,\frac{p^t(x^t-x^v)}{\sum_wp^wx^w}
\]
линейна: коэффициенты суть данные.

\textbf{Ошибки измерения.} Расходы: $p^tx^t-r_t\ge p^tx^v$, цель $\min\sum r_t/T$. Цены: нормировка $\sum_kp^t_k=1$, искомые $\tilde p^t$, ограничения $\tilde p^tx^t-\tilde p^tx^v\le\alpha U_{t,v}$ и $\alpha(U_{v,t}-1)\le\tilde p^tx^v-\tilde p^tx^t$, цель $\min\frac1T\sum_t\norm{p^t-\tilde p^t}_1$, линеаризованная через $\Delta^t_k\ge\pm(\tilde p^t_k-p^t_k)$. Все ограничения \emph{линейны} по $\tilde p^t$, потому что количества $x^t$ остаются данными. Это принципиальное отличие от нашей задачи о возмущении: у нас в ограничении свидетеля $\ip{\xi}{q_s-q_t}\ge1$ обе стороны~--- переменные, и член билинеен.

Размеры: $T$ от 22 до 79 наблюдений на субъекта, Gurobi, все индексы считаются за доли секунды, несмотря на NP-трудность.

\section{Стохастическая теория: популяция}

\subsection{Модель случайной полезности и аксиома ARSP}

Изложение по Stoye (2019), который даёт теорему McFadden--Richter в векторной форме.

\textbf{Примитивы.} Конечная вселенная альтернатив $X=(x_1,\dots,x_K)$, список задач выбора $\mathcal C=(C_1,\dots,C_J)$, $C_j\subset X$. Вероятности выбора $\pi(C_j)\in\Delta^{|C_j|-1}$ складываются в вектор $\Pi\in[0,1]^I$, $I=\sum_j|C_j|$, сумма координат $J$. Детерминированная функция выбора~--- бинарный вектор $R\in\set{0,1}^I$ с ровно одной единицей на задачу. $\mathcal R\subset\set{0,1}^I$~--- конечное множество рационализуемых функций выбора, в каком угодно смысле: максимизация полезности, GARP, IIA. \textit{Модель случайной полезности}: $\Pi$ порождён распределением на $\mathcal R$. Множество таких $\Pi$ есть выпуклая оболочка $\co(\mathcal R)$. \textit{Испытание}~--- вектор $T\in\set{0,1}^I$, единицы которого лежат внутри одной задачи выбора; $\mathcal T$~--- все испытания.

\textbf{ARSP.} Для любой конечной последовательности испытаний $(T_1,\dots,T_M)$
\[
    \sum_{m=1}^MT_m\Pi\ \le\ \max_{R\in\mathcal R}\sum_{m=1}^MT_mR.
\]

\textbf{Теорема (McFadden--Richter).} $\Pi$ стохастически рационализуем тогда и только тогда, когда выполнен ARSP.

\textbf{Доказательство в один абзац.} Последовательность входит только через $\sum_mT_m$, целочисленный неотрицательный вектор, и любой такой вектор представим. Условие продолжается на рациональные, затем на вещественные $T\ge0$ по непрерывности. Знак снимается трюком с нормировкой: $\mathbf 1\Pi=\mathbf 1R=J$, поэтому $T$ можно заменить на $T+\norm{T}_\infty\mathbf 1\ge0$. Остаётся: $T\Pi\le\max_RTR$ для всех $T\in\R^I$, то есть $\Pi$ не отделяется от $\co(\mathcal R)$ гиперплоскостью, то есть $\Pi\in\co(\mathcal R)$. Существенных испытаний конечное число: они отвечают градиентам фасет многогранника $\co(\mathcal R)$. Обобщение на многозначный выбор (GARSP) получается заменой $X$ на $2^X$, причём аксиома McFadden (2005) в этом случае оказывается слишком слабой.

Сравните с нашей сигнатурной леммой: там то же утверждение «$y\in\cone\set{\mathbf 1_S}$», а опровержение~--- вектор Фаркаша $\lambda$, который и есть «испытание» в смысле ARSP.

\subsection{Kitamura--Stoye: от непрерывного спроса к конусу}

\textbf{Данные.} $K$ благ, $J$ бюджетов $B_j=\set{y\in\Rp[K]\mid p_j'y=1}$, доход нормирован. Повторные срезы: по каждому бюджету наблюдается распределение выбора. Ограничения на агента: монотонность и SARP, то есть максимизация некоторой неубывающей полезности. Гетерогенность полезностей не ограничена ничем.

\textbf{Patches.} Для бюджета $B_j$ берётся самое грубое разбиение $\set{x_{1|j},\dots,x_{I_j|j}}$, при котором никакая другая бюджетная плоскость не пересекает внутренности элемента: для любых $k$ и $y_1,y_2\in x_{i|j}$ выполнено $(p_ky_1-1)(p_ky_2-1)\ge0$. Koida--Shirai записывают то же как класс эквивалентности по знакам: $y',y''$ в одном patch, если $\operatorname{sgn}(p_{j'}y'-1)=\operatorname{sgn}(p_{j'}y''-1)$ при всех $j'\ne j$. Patch кодируется строкой $X_i\in\set{-1,0,+1}^J$ («ниже, на, выше» каждого бюджета), существование patch с данным кодом проверяется совместностью системы $p_ky\lessgtr1$, $y\ge0$. $I=\sum_jI_j$~--- общее число patches. Выбор на пересечениях бюджетов имеет вероятность нуль и исключается.

Ключевое наблюдение: монотонность загоняет выбор на бюджетную плоскость, а SARP зависит только от того, выше или ниже каких бюджетов лежит выбор. Поэтому два агента, выбирающие из одних и тех же patches, либо оба рациональны, либо оба нет; распределение спроса важно только \emph{между} patches, не внутри. Непрерывная задача без потери информации становится дискретной с $I$ объектами.

\textbf{Типы и матрица $A$.} Тип~--- $a\in\set{0,1}^I$ с одной единицей на бюджет. Выбор $x_i$ из $B_j$ выявляет $x_i\succ x_m$ для всех $x_m$ с $X_{mj}\in\set{-1,0}$: на $B_j$ или ниже. Тип \textit{допустим}, если транзитивное замыкание ациклично (Флойд--Уоршелл). Допустимые типы $a_1,\dots,a_H$~--- столбцы $A\in\set{0,1}^{I\times H}$.

\textbf{Гипотезы.}
\begin{align*}
    (\mathrm H_A)&\quad \exists\nu\in\Delta^{H-1}:\ A\nu=\pi;\\
    (\mathrm H_B)&\quad \exists\nu\ge0:\ A\nu=\pi;\\
    (\mathrm H_C)&\quad \min_{\eta\in C}(\pi-\eta)'\Omega(\pi-\eta)=0,\qquad C:=\set{A\nu\mid\nu\ge0}.
\end{align*}
Эквивалентность $\mathrm H_A\Leftrightarrow\mathrm H_B$: $\mathbf 1_I'A=(J,\dots,J)$ и $\mathbf 1_I'\pi=J$, поэтому из $A\nu=\pi$ следует $\mathbf 1'\nu=1$ автоматически. Авторы прямо говорят, что с конусом $C$ удобнее, чем с многогранником. McFadden--Richter формулировали то же как ЛП: $\min_{\nu\ge0,s\ge0}\mathbf 1's$ при $A\nu+s\ge\pi$, $\mathbf 1'\nu\le1$.

\textbf{Статистика.} $J_N=N\min_{\nu\ge0}(\hat\pi-A\nu)'\Omega(\hat\pi-A\nu)$, где $\hat\pi$~--- выборочные частоты. Проекция $\hat\eta=A\hat\nu$ единственна, $\hat\nu$~--- нет. $J_N=0$ тогда и только тогда, когда $\hat\pi$ рационализуем.

\textbf{Вывод критических значений.} Предельное распределение $J_N$ разрывно зависит от того, сколько неравенств активно, наивный бутстрэп невалиден. Три подхода: регуляризация (субсэмплинг), отбор неравенств (требует $\mathcal H$-представления конуса, то есть исключения Фурье--Моцкина, что невозможно при большом $H$) и «подтягивание»: заменить $\nu\ge0$ на $\nu\ge\tau_N\mathbf 1_H$, $\tau_N\downarrow0$, $\sqrt N\tau_N\uparrow\infty$,
\[
    J_N(\tau_N)=\min_{\nu-\tau_N\mathbf 1\ge0}N(\hat\pi-A\nu)'\Omega(\hat\pi-A\nu),
\]
бутстрэп-выборки центрировать вокруг подтянутой проекции $\hat\eta_{\tau_N}$: $\hat\pi^{*(r)}_{\tau_N}=\hat\pi^{*(r)}-\hat\pi+\hat\eta_{\tau_N}$. Теорема~5.1: квантиль распределения $\tilde J_N(\tau_N)$ даёт равномерно валидный критический уровень. Практический выбор $\tau_N=\sqrt{\log\underline N/\underline N}$, $\underline N=\min_jN_j$.

\textbf{Вычисления.} Три способа построить $A$: грубая сила ($\prod_jI_j$ кандидатов), «crawling»~--- обход дерева по бюджетам с проверкой цикла в каждом узле и отсечением ветви, и предложение~4.1: если $B_1,\dots,B_M$ не пересекают $B_J$, то допустимые типы на $(B_1,\dots,B_J)$ суть произведения допустимых типов на $(B_1,\dots,B_{J-1})$ и на $(B_{M+1},\dots,B_J)$. В приложении (UK FES, 1975--1999, окна по 8 лет, 3 агрегированных блага, по одному бюджету в год на медианном расходе): $I$ от 14 до 67, $H$ от 21 до 177352, $J_N=0$ ровно в одном окне, но нарушения незначимы ($p$-значения от 9\% до 100\%). Построение $A$~--- от долей секунды до часов, $J_N$~--- до минуты.

\textbf{Пример $J=2$.} Два пересекающихся бюджета, по два patch:
\[
    X=\begin{pmatrix}0&-1\\0&1\\-1&0\\1&0\end{pmatrix},\qquad
    A=\begin{pmatrix}1&0&0\\0&1&1\\0&0&1\\1&1&0\end{pmatrix}.
\]
Тип $(1,0,1,0)'$ нарушает слабую аксиому и в $A$ не входит. При $\Omega=I$ получается $J_N=N\max\set{\hat\pi_1+\hat\pi_3-1,0}^2$: всё содержание рациональности~--- одно неравенство $\pi_1+\pi_3\le1$.

\textbf{Частичная идентификация.} Идентифицированное множество
\[
    \mathbf H_\nu=\set{\nu\in\Delta^{H-1}\mid A\nu=\pi}
\]
~--- многогранник высокой размерности; при $\pi$ без нулей и $I\ge4$ он никогда не точка. Контрфактические оценки: для линейной $f(\nu)=e'\nu$ границы $\min/\max\,e'\nu$ при $A\nu=\hat\pi$, $\nu\ge0$~--- задачи ЛП; так ограничиваются вероятности patches на новом бюджете, ожидаемый спрос и функция распределения спроса (следствия 8.1--8.3).

\textbf{Замечание 8.3 статьи.} Для дискретного выбора, в частности бинарного, шаг дискретизации не нужен, и метод применим напрямую. Это тот случай, в котором живёт наша модель: решение агента при каждой цене бинарное.

\subsection{Генерация столбцов: Smeulders}

Проблема: $|\mathcal R|$ растёт экспоненциально по $T$ (таблица~1: 3 блага, 20 периодов~--- до $3\cdot10^{22}$ типов), а QP нужно решать $1+M$ раз, по разу на бутстрэп-выборку.

\textbf{Главная задача.} QP (4)--(6): $\min N\sum_{t,i}s_{t,i}^2$ при $\sum_{r\in\mathcal R_{t,i}}p_r+s_{t,i}=\hat\pi_{t,i}$, $p_r\ge0$. Ограниченная главная задача~--- то же над подмножеством $\bar{\mathcal R}$, решение $(\bar p^*,\bar s^*)$, проекция $\bar v^*=\sum\bar p^*_ra_r$.

\textbf{Критерий остановки (теоремы 2--3).} $v^*$~--- проекция $\hat\pi$ на конус тогда и только тогда, когда $(\hat\pi-v^*)\cdot(a_r-v^*)\le0$ для всех образующих $a_r$. Поэтому ограниченное решение оптимально, если
\[
    \max_{r\in\mathcal R}\bar s^*\cdot a_r\ \le\ \bar s^*\cdot\bar v^*,\tag{11}
\]
иначе столбец-нарушитель добавляется. Это стандартная проверка приведённой стоимости: $\bar s^*$ играет роль двойственных цен. Достаточно любого столбца выше порога, оптимум подзадачи не нужен.

\textbf{Подзадача ценообразования (GAPP-версия, (19)--(25)).} Данные: $s_{t,i}$ из $\bar s^*$; $X_{t,i,j}\in\set{0,1}$, равное $1$, если выбор patch $x_{t,i}$ влечёт $p_j\succ_pp_t$. Переменные: $\alpha_{t,i}\in\set{0,1}$ (patch выбран), $\rho_{t,j}\in\set{0,1}$ (период $t$ выявленно предпочтён $j$).
\begin{empheq}[left=\empheqlbrace]{align*}
    &\sum_{t,i}s_{t,i}\alpha_{t,i}\ \to\ \max,\\
    &\sum_i\alpha_{t,i}=1, & &\forall t,\\
    &\sum_i\alpha_{t,i}X_{t,i,j}-\rho_{j,t}\le0, & &\forall j,t,\\
    &\rho_{j,t}+\rho_{t,k}-\rho_{j,k}\le1, & &\forall j,t,k,\\
    &\rho_{j,t}+\rho_{t,j}\le1, & &\forall j,t.
\end{empheq}
Третье ограничение заставляет $\rho$ содержать прямые отношения, четвёртое~--- транзитивность, пятое~--- антисимметрию. Вместе они исключают циклы \emph{без} перечисления циклов: это MILP задачи линейного упорядочения, а не генерация отсекающих плоскостей по циклам. Перед точным MILP пробуется эвристика «лучшая вставка»: допустимый тип отвечает линейному порядку периодов, и при заданном порядке лучший тип считается покоординатно по (26).

\textbf{Подтягивание под генерацию столбцов.} Ограничение $p_r\ge\tau_N/|\mathcal R|$ для всех $r$ несовместимо с частичным набором столбцов. Лемма~5 сдвигает переменные, но правая часть всё ещё требует всего $\mathcal R$. Лемма~6: достаточно подмножества $\mathcal R'$, которое «задевает» каждую активную гиперплоскость $\mathcal H$-представления; тогда подтянутый конус строго внутри исходного. $\mathcal R'$ из 1000 типов строится полуслучайно (приложение~B). Нижняя и верхняя оценки позволяют прерывать генерацию столбцов на бутстрэп-выборке, как только ясен знак $J_N^{*(r)}-J_N$.

\textbf{Результаты.} CPLEX, UK FES: 10 периодов~--- 27 с против 103 с точным ценообразованием, 15 периодов~--- 565 с против недоступного, 20 периодов~--- около 7 часов на 1000 выборок; DKQS до этого ограничивались 8 периодами.

\subsection{Двойственная характеристика: Koida--Shirai}

Та же задача, но конус описан неравенствами, а не образующими.

\textbf{Конструкция $\Xi$.} Для каждого подсемейства бюджетов $\mathcal J\subset[J]$, $|\mathcal J|\ge2$, patch $B_{ji}$ называется \textit{недоминируемым в $\mathcal J$}, если $j\in\mathcal J$ и $p_{j'}y-1>0$ для всех $y\in B_{ji}$, $j'\in\mathcal J\setminus\set j$: ни один другой бюджет семейства не выявляет предпочтение над ним. Строка $\xi^{\mathcal J}\in\set{0,1}^I$~--- индикатор недоминируемых patches. $\Xi$~--- матрица из $2^J-J-1$ строк.

\textbf{Лемма 1.} Тип $a$ допустим тогда и только тогда, когда $\Xi a\ge\mathbf 1$: цикл на $\mathcal J$ делает все выбранные в $\mathcal J$ patches доминируемыми, и строка даёт нуль.

\textbf{Теорема 1.} $\pi$ рационализуем тогда и только тогда, когда $\Xi\pi\ge\mathbf 1$.

Нетривиальный шаг: многогранник $\set{\pi\mid\Xi\pi\ge\mathbf 1}$ целочислен, то есть равен выпуклой оболочке своих целых точек. Доказывается через ранг Хватала, равный нулю: для любого целого $u\Xi$ сумма $u\cdot\mathbf 1$ целая, что показывается упорядочением бюджетов по общему направлению $d$ и разбиением $\mathcal J_{\!*}$ на блоки. Полная унимодулярность не используется, $\Xi$ нарушает стандартные достаточные условия.

\textbf{Пример $n=2$, $J=3$}, по три patch на бюджет:
\[
    \Xi=\begin{pmatrix}1&0&0&0&1&0&0&0&1\\1&0&0&0&1&1&0&0&0\\0&0&0&1&1&0&0&0&1\\1&1&0&0&0&0&0&1&1\end{pmatrix}\!,
\]
строки отвечают $\set{1,2,3},\set{1,2},\set{2,3},\set{1,3}$. Для $\pi=(\frac1{10},\frac8{10},\frac1{10};\frac4{10},\frac2{10},\frac4{10};\frac1{10},\frac8{10},\frac1{10})$ получается $\Xi\pi=(\frac25,\frac7{10},\frac7{10},\frac95)$: нарушены строки $\set{1,2}$ и $\set{2,3}$, и в любом представлении циклы на этих парах имеют положительный вес. Пример KS с $n=3$, $J=3$ даёт $\pi$, удовлетворяющий всем парным строкам (WARP), но с $\xi^{\set{1,2,3}}\pi=0$: строки размера $\ge3$ несут информацию сверх WARP.

\textbf{Теорема 2 и индекс.} Двойственная ЛП $\min\xi\cdot\pi$ при $\xi A\ge c$ достигает оптимума на одной из явных строк: $\min_{\mathcal J}\xi^{\mathcal J}\cdot\pi=\max\set{\sum_{a\in\mathcal A^*}\tau_a\mid A\tau=\pi,\tau\ge0}$~--- наибольший вес, который можно положить на рациональные типы. Это непрерывный индекс рациональности для популяции.

Для нас это прямой аналог вопроса о фасетах конуса $\cone\set{\mathbf 1_S\mid S\in\Sp(P)}$ и предложений 5.2--5.4 Шананина--Молчанова о дискретной выпуклости: там грани с нормалями из $\set{-1,0,1}$, здесь строки из $\set{0,1}$ с правой частью $1$.

\subsection{Стохастическая версия DKQS: RAUM}

\textbf{Определение.} Повторные срезы $\set{(p^t,\hat\pi^t)}$. Данные рационализуемы RAUM, если существует вероятностное пространство и $\chi:\Omega\to(\Rp[L])^T$ такие, что почти все траектории $\set{(p^t,\chi^t(\omega))}$ удовлетворяют GAPP и $\hat\pi^t(Y)=\mu\set{\chi^t\in Y}$.

\textbf{Сведение.} По предложению~1 GAPP равносилен GARP на изорасходных данных, поэтому спрос проецируется радиально на плоскость $B^t=\set{x\mid p^tx=1}$, и далее работает машина KS: patches~--- ячейки разбиения плоскостей, типы~--- классы GARP-эквивалентности, матрица $A$. \textbf{Теорема 3.} Данные рационализуемы RAUM тогда и только тогда, когда $Av=\pi$ при некотором $v\in\Delta^{|\mathcal A|-1}$. Преимущество перед прямым RUM-тестом: не нужно оценивать распределение спроса при фиксированном доходе.

\textbf{Благосостояние.} Доля типов, которым выгодно изменение цен $p^{t'}\to\hat p$, равна $\theta=\rho\cdot v$ с $\rho=\mathbf 1_{\hat p\succeq_p^*p^{t'}}$; точные границы $\min/\max\rho\cdot v$ при $Av=\pi$, $v\ge0$~--- ЛП. Множество $\mathcal S(\theta)=\set{Av\mid\rho\cdot v=\theta,v\in\Delta}$~--- многогранник, не конус, поэтому бутстрэп KS не работает; вводится «подтягивание, зависящее от ограничения» (теорема~4). Приложения: Progresa (83\% домохозяйств удовлетворяют GAPP), UK FES и канадский SHS (RAUM не отвергается, доверительные интервалы для долей шириной менее $0{,}1$--$0{,}15$).

\subsection{Ошибки измерения в популяции: Aguiar--Kashaev}

Это другой способ соединить гетерогенность и возмущение данных: не минимизировать возмущение, а допустить его как латентную переменную с моментным ограничением.

\textbf{Данные.} Панель $n$ потребителей за $T$ периодов, $(\rho_t,c_t)$, эффективные цены $\rho_t$ (статика: $p_t$; динамика с экспоненциальным дисконтированием: $p_t/\prod_{j\le t}(1+r_j)$).

\textbf{Модель.} Истинные $(\rho_t^*,c_t^*)$ латентны, наблюдаемые отличаются ошибкой $w_t$ без предположений об аддитивности или независимости. Рационализуемость в форме условий первого порядка $\delta_t\nabla u(c_t^*)\le\lambda_t\rho_t^*$, после исключения $u$~--- неравенства Африата со случайными числами полезности:
\[
    v_t-v_s\ \ge\ \frac{\lambda_t}{\delta_t}\rho_t^{*\prime}(c_t^*-c_s^*)\quad\text{почти наверное}.
\]
Без ограничения на $w$ любые данные совместимы с любой моделью. Единственное ограничение~--- центрирующий момент $\mathbb E[g_M(x,e)]=0$: средняя бюджетная нейтральность $\mathbb E[\rho_t'w^c_t]=0$ для обследований, $\mathbb E[w^c_t]=0$ для «дрожащей руки» в экспериментах, $\mathbb E[\rho_t]=\mathbb E[\rho_t^*]$ для ошибки восприятия цен.

\textbf{Проверка.} Неравенства записываются индикаторными моментами $g_I$. Теорема~2: рационализуемость равносильна $\inf_\mu\norm{\mathbb E_{\mu\times\pi_0}g}=0$ по всем условным законам латентных переменных~--- бесконечномерная задача. ELVIS (Schennach 2014) сводит её к конечномерной через максимум энтропии: $h(x;\gamma)$~--- среднее $g$ под экспоненциально наклонённой мерой, и рационализуемость равносильна $\inf_\gamma\norm{\mathbb E h(x;\gamma)}=0$ (теорема~3). Теорема~4: двойственные к неравенствам уходят в $+\infty$, остаётся минимизация по $\gamma_M\in\R^q$, а неравенства превращаются в ограничение носителя, с которого сэмплируют отбраковкой или MCMC. Статистика $\mathrm{TS}_n=n\inf_\gamma\hat h_M'\hat\Omega^-\hat h_M$, предел $\chi^2_q$.

\textbf{Результаты.} Детерминированный тест Браунинга отвергает экспоненциальное дисконтирование у 84--90\% домохозяйств; с ошибкой измерения~--- не отвергается у одиночек, отвергается у пар. В эксперименте Ahn et al. детерминированный тест проходят 13\% субъектов; с ошибкой восприятия цен модель не отвергается, с ошибкой в количествах~--- отвергается.

\section{Сравнение с нашей постановкой}

\begin{center}
\renewcommand{\arraystretch}{1.25}
\begin{tabular}{>{\raggedright\arraybackslash}p{0.47\textwidth}>{\raggedright\arraybackslash}p{0.47\textwidth}}
\hline
Ветка выявленных предпочтений & Заметки о разрешимости\\
\hline
$J$ бюджетов~--- гиперплоскости в пространстве благ & $T$ цен~--- поверхности $\set{h(p_t\circ x)=1}$ в пространстве агентов\\
Patches: классы по знакам $\operatorname{sgn}(p_{j'}y-1)$ на самих гиперплоскостях & Камеры $R_S$: классы по знакам $h(p_t\circ x)-1$, полномерные\\
Тип: один patch на бюджет & Спектр $S\subset[T]$\\
Допустимость: SARP, Флойд--Уоршелл; в MILP~--- переменные порядка с транзитивностью и антисимметрией & Достижимость: $t\in S$ не покрывает $[T]\setminus S$, ЛП; в MILP~--- свидетели $\xi$\\
Существование patch: ЛП на $p_ky\lessgtr1$ & Покрытие: ЛП на веса $w$\\
$\pi=A\nu$, $\nu\ge0$ & $y\in\cone\set{\mathbf 1_S\mid S\in\Sp(P)}$\\
ARSP: $T\Pi\le\max_RTR$; двойственно $\Xi\pi\ge\mathbf 1$ & Опровержение Фаркаша $\lambda$; предл. 5.2 Шананина\\
QP-расстояние, бутстрэп с подтягиванием & Точное ЛП, сертификат\\
Генерация столбцов, подзадача~--- MILP упорядочения & Перебор $2^T$ с отсечением\\
Предл.~4.1: непересекающиеся бюджеты~--- произведение типов & Доминирующие группы цен~--- произведение спектров\\
$\pi_1+\pi_3\le1$ при $J=2$ & $y_2\le y_1+y_3$ при $T=3$\\
Пересечения бюджетов~--- событие нулевой вероятности & Общее положение не требуется, доказано\\
Гетерогенность максимальная: своя полезность & Гетерогенность минимальная: общая $h$, свой $x$\\
Возмущение цен (APE): линейно, количества фиксированы & Возмущение цен: билинейно, свидетели переменные\\
Ошибка измерения как латентная переменная с моментом & Минимальное возмущение с глобальным оптимумом\\
\hline
\end{tabular}
\end{center}

\paragraph{Источник ограничений.} У KS ограничения на $\pi$ идут из рациональности каждого агента при разных бюджетах. У нас $h$ одна на всех, решение бинарное. KS в замечании~8.3 сами отмечают, что для дискретного выбора их метод применим без дискретизации; но при бинарном выборе без общей структуры допустимы все $2^T$ типов, и конус тривиален. Наш критерий~--- это в точности те ограничения, которые накладывает общая технология. Две задачи~--- противоположные концы одной конструкции.

\paragraph{Где живёт комбинаторика.} У них~--- в пространстве благ, через положение patch относительно гиперплоскостей и ацикличность. У нас~--- в пространстве цен, через $p_t\in\conv\set{p_s}+\Rp[d]$. DKQS переносят отношение в пространство цен, но типы у них по-прежнему классы GARP на изорасходных данных. Поэтому у нас цены можно сделать переменными; у них при возмущении цен (Demuynck--Rehbeck) количества остаются данными, и задача остаётся линейной.

\paragraph{Два способа возмущать.} Линия Варьяна и Demuynck--Rehbeck минимизирует возмущение одного агента при фиксированной комбинаторике отношений. Aguiar--Kashaev не минимизируют: ошибка~--- латентная переменная, ограниченная моментом, а тест~--- существование закона. Наша задача ближе к первой линии, но с геометрическим, а не порядковым условием допустимости типа, отсюда билинейность и необходимость глобального солвера.

\paragraph{Экономический смысл.} Однородная вогнутая $h(p)$~--- функция удельных издержек технологии с постоянной отдачей; $h(p_t\circ x)<1$~--- фирма с факторной эффективностью $x$ при ценах $p_t$ прибыльна. У них~--- потребители, бюджеты, спрос.

\section{Что стоит взять}

\begin{itemize}
    \item \emph{Генерация столбцов.} Главная задача та же: ЛП или QP над столбцами $\mathbf 1_S$. Подзадача ценообразования: найти достижимый $S$ с $\bar s\cdot\mathbf 1_S>$ порог. По лемме о свидетелях это MILP с бинарными $z_t$ и переменными $\xi_t$; при фиксированных ценах ограничение $\ip{\xi_t}{p_s-p_t}\ge1-M(1-z_t+z_s)$ линейно, так что подзадача~--- обычный MILP без билинейности. Эвристика «лучшая вставка» не переносится: у покрытия нет структуры линейного порядка.
    \item \emph{Двойственная форма.} Явная система неравенств для $\cone\set{\mathbf 1_S\mid S\in\Sp(P)}$ по образцу $\Xi\pi\ge\mathbf 1$, с проверкой целочисленности через ранг Хватала,~--- это вопрос о структуре опровержений $\lambda$ и о дискретной выпуклости Молчанова.
    \item \emph{Статистика.} Если $y$ наблюдается с шумом, статистика $J_N$ и бутстрэп с подтягиванием переносятся дословно, конус у нас тот же. Для частично идентифицированных функционалов от $\nu$ (у нас от масс $m_S$)~--- процедура DKQS.
    \item \emph{Что не брать.} SARP, порядковые переменные, изорасходную нормировку: в нашей модели нет отношения предпочтения.
\end{itemize}

\section{Атрибуция}

Принцип «смесь по ячейкам есть конус индикаторов» в нашей сигнатурной лемме~--- предложение 5.1 Шананина (1999) при фиксированной $h$ и, независимо, Block--Marschak, McFadden--Richter, Kitamura--Stoye. Постановка «ближайшие данные, при которых всё согласовано»~--- Varian (1985), MILP-техника с big-M и бинарными индикаторами~--- Demuynck--Rehbeck. Эти источники должны быть названы во введении. Характеристика достижимых спектров через покрытие, реализуемость всего $\Sp(P)$ одной $h$, отказ от общего положения, произвольное $d$, полный спектр на гиперболе и возмущение цен как билинейная MIQCP в перечисленных работах отсутствуют.

\section*{Что читать}

\begin{enumerate}
    \item Y.~Kitamura, J.~Stoye. Nonparametric analysis of random utility models. \emph{Econometrica} 86 (2018) 1883--1909. Рабочая версия 2012: \url{https://www.princeton.edu/~erp/NSF%20Conference/PDFs/Stoye%20nprum3_9_6.pdf}.
    \item J.~Stoye. Revealed stochastic preference: a one-paragraph proof and generalization. \emph{Economics Letters} (2019). \url{https://arxiv.org/abs/1810.10604}.
    \item B.~Smeulders. Column generation algorithms for nonparametric analysis of random utility models. \url{https://arxiv.org/abs/1812.01400}; журнальная версия: Smeulders, Cherchye, De Rock, \emph{Econometrica} 89 (2021) 437--455.
    \item K.~Koida, K.~Shirai. A dual approach to nonparametric characterization for random utility models. \url{https://arxiv.org/abs/2403.04328}.
    \item R.~Deb, Y.~Kitamura, J.~Quah, J.~Stoye. Revealed price preference: theory and empirical analysis. \emph{REStud} 90 (2023) 707--743. \url{https://kitamura.sites.yale.edu/sites/default/files/files/DKQS.pdf}.
    \item H.~Varian. Goodness-of-fit for revealed preference tests (1990). \url{https://econwpa.ub.uni-muenchen.de/econ-wp/em/papers/9401/9401001.pdf}. Его же: Non-parametric analysis of optimizing behavior with measurement error. \emph{J.~Econometrics} 30 (1985) 445--458.
    \item M.~Dean, D.~Martin. Measuring rationality with the minimum cost of revealed preference violations. \url{https://www.columbia.edu/~md3405/Working_Paper_11.pdf}.
    \item T.~Demuynck, J.~Rehbeck. Computing revealed preference goodness-of-fit measures with integer programming. \emph{Economic Theory} (2023); ECARES WP 2021-26: \url{https://dipot.ulb.ac.be/dspace/bitstream/2013/334880/3/2021-26-DEMUYNCK_REHBECK-computing.pdf}.
    \item V.~Aguiar, N.~Kashaev. Stochastic revealed preferences with measurement error. \emph{REStud} 88 (2021) 2042--2093. \url{https://arxiv.org/abs/1810.05287}.
\end{enumerate}

\end{document}
